\section{Limitations and Threats to Validity}
\label{sec:limits}

\emph{No physical measurement.} No FPGA board was available. Every time in
this paper is a model value; $\tau$, $T_{\mathrm{host}}$, the behaviour of
the USB--UART bridge, board-level reliability and power are unmeasured. A
claim that fusion or residency is faster on hardware requires measured
distributions of $\tau$ and call times on a programmed board, a stated
warm-up and repetition policy, and a timed CPU baseline for the same
kernels~\cite{hoefler-benchmarking}. The supplement fixes such a procedure
and the form of its report.

\emph{Model validity.} Eq.~\eqref{eq:time} assumes one turnaround for every
command; buffering may make $\tau$ depend on the command and reply size,
which would shift the break-even values. Counts come from a clockless model;
their transfer to the RTL rests on byte-exact replay of recorded transcripts,
not on an equivalence proof.

\emph{Workloads.} Five fixed kernels with fixed sizes and inputs are
studied; matrix zero skipping and resident byte stuffing make counts
input-dependent. The repeated identical dot product is an amortisation
microbenchmark. The geometry comparison holds problem sizes fixed but not
mappings optimised per shape.

\emph{Verification scope.} The proofs cover selected safety properties on
enumerated instances; completion of transactions is shown only under
explicit delivery and service obligations. The legacy protocol's recovery
from arbitrary byte insertion or loss is not proved, and the formal
properties do not cover the C library. The link between proved RTL
registers and constrained netlist cells is by name (Section~\ref{sec:chain}).

\emph{Implementation evidence.} Each geometry was placed and routed once,
on an earlier snapshot whose wrapper differs from the current one only by
wiring; only the $4\times4$ point was re-implemented on the current
revision. The period search is heuristic and was not repeated; its candidate
was re-checked on the current revision.
Resources and frequencies are specific to one device, tool version and
strategy, and are not comparable with other work.

\emph{Same-author references.} The device model, RTL, test benches and
properties were written by the same author; independent scalar references
and mutation tests reduce, but do not remove, the risk of a shared
misunderstanding.
